%% kapa: exam=1 round=30 year=2019 period=1 form=A subject=economics subject_ko=경제학원론 seq=37 number=77
%% source: Q-Net 감정평가사 기출문제 (cst00309) artlSeq=5205205 file=extracted/1008/5205205/2019년 제30회 감정평가사 1차시험 문제지/2019년 제30회 감정평가사 1차시험 문제지/2019년도 제30회 감정평가사 1차 1교시 A형.hwp
%% conversion: kordoc 4.24.51 -> tools/gen1.py (auto); review=proofread-sonnet; flags=none
\begin{question}{77}
통화공급 과정에 관한 설명으로 옳은 것을 모두 고른 것은?
\begin{framedtext}
ㄱ. 100 \% 지급준비제도가 실행될 경우, 민간이 현금통화비율을 높이면 통화승수는 감소한다. \par
ㄴ. 민간이 현금은 보유하지 않고 예금만 보유할 경우, 예금은행의 지급준비율이 높아지면 통화승수는 감소한다. \par
ㄷ. 중앙은행이 민간이 보유한 국채를 매입하면 통화승수는 증가한다.
\end{framedtext}
\begin{choices}
  \item ㄱ
  \item ㄴ
  \item ㄱ, ㄴ
  \item ㄱ, ㄷ
  \item ㄴ, ㄷ
\end{choices}
\end{question}
\answer{2}
